\clearpage
\section{SOLUSI DAN TARGET LUARAN}

\subsection{Solusi}

Menanggapi permasalahan yang telah diidentifikasi dalam sub bab 1.2, solusi utama yang ditawarkan dalam kegiatan Pengabdian kepada Masyarakat ini adalah penyelenggaraan Workshop \textit{Machine Learning} untuk Prediksi Awal Kesehatan Tanaman melalui Analisis Citra Daun bagi siswa Jurusan Rekayasa Perangkat Lunak SMKN 2 Cimahi. Kegiatan ini dirancang sebagai transfer pengetahuan dan keterampilan praktis agar siswa tidak hanya memahami konsep dasar kecerdasan buatan, tetapi juga mampu membangun proyek sederhana berbasis \textit{Machine Learning} dan \textit{Computer Vision}.

Workshop menggunakan pendekatan \textit{Project-Based Learning} yang menempatkan siswa sebagai peserta aktif dalam proses belajar. Pendekatan ini sesuai dengan karakter pendidikan vokasional yang menekankan keterampilan praktik, penyelesaian masalah, kolaborasi, dan pembuatan produk nyata. Siswa diarahkan membangun alur kerja mulai dari pengenalan \textit{dataset}, praproses gambar, pelatihan model klasifikasi, evaluasi hasil prediksi, hingga pembuatan aplikasi sederhana yang dapat menerima input foto daun dan menampilkan hasil prediksi, sehingga memperoleh pengalaman membangun \textit{pipeline Machine Learning} secara \textit{end-to-end}.

Studi kasus klasifikasi citra daun dipilih berdasarkan beberapa pertimbangan: \textit{dataset} tersedia secara publik, data berbentuk gambar lebih mudah divisualisasikan, topik kesehatan tanaman relevan dengan isu ketahanan pangan dan penerapan kecerdasan buatan dalam kehidupan sehari-hari, serta tidak memerlukan perangkat keras khusus sehingga cukup menggunakan laptop dan koneksi internet. Workshop memanfaatkan perangkat lunak gratis berbasis browser, yaitu \textit{Google Colab} untuk pelatihan model, Kaggle sebagai sumber dataset, dan \textit{Streamlit} untuk pengembangan aplikasi prediksi sederhana. Modul pelatihan dan \textit{notebook} praktik disiapkan secara terstruktur agar siswa dapat mengulang materi secara mandiri setelah kegiatan selesai.

Hasil prediksi sistem yang dibangun dijelaskan kepada siswa sebagai alat \textit{screening} awal, bukan diagnosis final. Melalui penegasan ini, siswa tidak hanya belajar membangun model, tetapi juga memahami keterbatasan model, pentingnya kualitas data, dan perlunya interpretasi hasil secara kritis, sehingga kegiatan ini turut membangun kemampuan berpikir ilmiah dan literasi kecerdasan buatan. Solusi ini diharapkan menjawab permasalahan utama mitra terkait keterbatasan akses pembelajaran \textit{Machine Learning} yang praktis, minimnya pengalaman proyek berbasis kecerdasan buatan, serta belum tersedianya portofolio AI/ML siswa, sekaligus menghasilkan modul pelatihan yang dapat digunakan kembali oleh sekolah mitra.

\subsection{Target Luaran}

Target luaran dalam kegiatan Pengabdian kepada Masyarakat ini disusun untuk menjawab permasalahan mitra secara terukur dan berdampak. Adapun target luaran dari program ini meliputi:

\subsubsection{Luaran Produk/Jasa}
\begin{enumerate}[label=\alph*), leftmargin=*]
    \item Peningkatan kompetensi siswa melalui pelatihan praktis \textit{Machine Learning} dan \textit{Computer Vision} berbasis studi kasus klasifikasi citra daun tanaman.
    \item Modul pelatihan \textit{Machine Learning} dan \textit{notebook Google Colab} beserta aplikasi prediksi sederhana berbasis \textit{Streamlit} yang dapat menjadi bahan ajar tambahan di sekolah mitra.
\end{enumerate}

\subsubsection{Luaran Artikel Ilmiah dan Dokumentasi}-`'
\begin{enumerate}[label=\alph*), leftmargin=*]
    \item Laporan Kegiatan dalam bentuk hard file dan soft file yang memuat seluruh proses pelaksanaan kegiatan.
    \item Kuesioner Evaluasi Pembelajaran untuk mengukur efektivitas pelatihan terhadap peningkatan kompetensi siswa.
    \item Dokumentasi hasil mini project siswa berupa proyek klasifikasi citra daun yang dikembangkan selama workshop.
    \item Poster kegiatan PKM dalam format A2 sebagai media informasi dan diseminasi pada tahap monitoring dan evaluasi.
    \item Publikasi Ilmiah berupa satu artikel ilmiah yang ditargetkan untuk diterbitkan di jurnal pengabdian masyarakat bereputasi (minimal memiliki ISSN dan terindeks dalam portal seperti Sinta).
\end{enumerate}
